\section{Extending (Almost) All That DuckDB Has to Offer}
\label{sec:strategy}
\label{sec:hooks}
% Shrunk 9 Nov 2026 to the interface table and one paragraph; the threshold discussion lives in 3.4, 4.4 and 4.7.

\rev{Table~\ref{tab:hooks} lists the \duckdb interfaces the operator and the index use, what each one gave us and
what it lacked, with the sections where a lack cost a mechanism; it is the list the introduction promised. Everything
runs through these interfaces without a change to \duckdb, except that the visibility check calls into two classes
outside the extension API, the transaction and the table storage. Three additions would make an index and an operator
like ours simpler. A transaction or commit identifier passed with each index callback would let the index decide
visibility itself instead of asking \duckdb once per change group (Section~\ref{sec:index-read}). A visibility check
that stays valid after version garbage collection would let a fold decide at any time what an epoch accounts for,
instead of while its transaction is live (Section~\ref{sec:index-fold}). Tasks for index work and task groups for
operators would let a fold run in the background and a traversal schedule its levels without a chain of events
(Sections~\ref{sec:phases} and~\ref{sec:index-fold}). Two entries of the table are not gaps but promises the interface
does not make, that row ids stay stable while an index exists and that \texttt{CanFetch} answers the visibility
question of snapshot isolation (Section~\ref{sec:index-correctness}); an extension that relies on them, as ours does,
should say so.}

\begin{table}[t]\centering\scriptsize\setlength{\tabcolsep}{3pt}
  \caption{The \duckdb extension interfaces the operator and the index use, what each provided, and what it lacked
  for this work, with the sections where the lack cost a mechanism. The graph never leaves the relational tables: no
  storage extension, custom catalog, replacement scan or custom type is involved.}
  \label{tab:hooks}
  \begin{tabular}{p{1.9cm}p{2.9cm}p{2.9cm}}\toprule
  interface & used for & what it lacked \\ \midrule
  parser extension & \sqlpgq syntax: \texttt{MATCH}, \texttt{CREATE PROPERTY GRAPH} & -- \\
  \addlinespace[2pt]
  table function with bind replacement & \texttt{MATCH} expanding to the rewritten SQL & -- \\
  \addlinespace[2pt]
  scalar functions & the marker that carries the path pattern to the planner; index statistics; strategy introspection & a marker must be volatile to survive folding, and must not reach execution (\S\ref{sec:planner}) \\
  \addlinespace[2pt]
  optimizer extension & replacing the marker and its two inputs by the logical operator & runs before or after \duckdb's optimizer, not between its rules; the rewrite must find its inputs on the binder's plan shape (\S\ref{sec:planner}) \\
  \addlinespace[2pt]
  logical and physical operator & sink for pairs and edges, source for results, own events and tasks & a cardinality, but no cost, and join ordering treats the operator as opaque (\S\ref{sec:choice}); no task groups in the scheduler (\S\ref{sec:phases}) \\
  \addlinespace[2pt]
  index type & build, append and delete callbacks; serialization into the block manager; WAL replay & rows but not transactions; no visibility at past points in time (\S\ref{sec:index-interface}, \S\ref{sec:index-fold}); stable row ids and \texttt{CanFetch} are implementation behaviour, not guarantees (\S\ref{sec:index-correctness}); no scan interface, so the operator finds the index in the catalog by name; no tasks for index work, so its build, fold and load start threads of their own (\S\ref{sec:index-fold}) \\
  \addlinespace[2pt]
  client context state & per-connection registry of property graphs and operator state & -- \\
  \addlinespace[2pt]
  settings & operator and index options, per connection or global & index options are fixed at \texttt{CREATE INDEX}, so a policy such as the fold ratio is a setting and not a property of the index (\S\ref{sec:index-fold}) \\
  \addlinespace[2pt]
  buffer manager & the multi-source working sets, the build's staged edges and the persisted index segments & -- \\
  \bottomrule\end{tabular}
\end{table}
